\chapter{Resultados}
\label{cap:resultados}

% ============================================================================
% Este capítulo se redactará en la fase de evaluación (F7).
% El protocolo de evaluación debe definirse ANTES de ejecutar los experimentos,
% para evitar el sesgo de fijar las métricas una vez conocidos los resultados.
% ============================================================================

\section{Protocolo de evaluación}
% TODO: conjunto de datos de evaluación, verdad conocida, líneas base de comparación,
% número de repeticiones por experimento.

\section{Métricas de detección}
% TODO: precisión, exhaustividad y F1 por técnica; tasa de falsos positivos sobre la
% porción benigna; cobertura ATT&CK.

\section{Métricas de correlación multi-etapa}
% TODO: campañas correctamente reconstruidas; completitud de etapas; ratio de reducción
% de alertas frente a la línea base sin correlación.

\section{Métricas de priorización}
% TODO: presencia del incidente real en el top-K; correlación entre la puntuación
% asignada y la gravedad anotada.

\section{Métricas de razonamiento}
% TODO: coincidencia entre la ruta derivada y la ruta real del plan de emulación.

\section{Rendimiento}
% TODO: rendimiento por etapa, latencia extremo a extremo, consumo de recursos,
% comportamiento bajo carga acelerada.

\section{Discusión}
% TODO: análisis de los resultados, casos en que el sistema falla y por qué,
% limitaciones del enfoque de evaluación.
